% GENERATED FILE. Edit canonical lessons, then run npm run generate:books.
% canonical-source-hash: fnv1a64:b836cf1d7905e0e5
% canonical-lessons: LA-C377-fullo, LA-C377-aurifex, LA-C377-lignarius, LA-C377-caementarius, LA-C377-statuarius

\chapter{Trades of the Town}
\label{ch:la-trades-of-the-town}
\hlchaptermodality{377}

\begin{chapteropening}
\textbf{By the end of this chapter:} \emph{I can name a fuller, a goldsmith, a carpenter, a mason and a sculptor.}

Complete the last lesson of chapter 377: I can name a fuller, a goldsmith, a carpenter, a mason and a sculptor.
\end{chapteropening}

\section[fullō, fullōnis]{fullō, fullōnis — a fuller, a cleaner of cloth}

\label{lesson:LA-C377-fullo}



\begin{quote}
Before the new one: say the Latin for a labourer, then the Latin for to work.
\end{quote}

\subsection*{You'll want to know: fullō, fullōnis}
\textbf{fullō, fullōnis} — \textquotedblleft{}a fuller, a cleaner of cloth\textquotedblright{}.

\begin{quote}
Fullō vestēs sordidās lavat.
\end{quote}

\begin{quote}
\emph{The fuller washes dirty clothes.}
\end{quote}

\emph{fullō, fullōnis} m. The fuller cleaned and finished woollen cloth by treading it in tubs; his shop was the \emph{fullōnica}.

\begin{cousinweb}
The origin of \emph{fullō} is not settled: Lewis \& Short give none, and EtymDB has no parent for it.
\end{cousinweb}

\begin{grammarlens}[title={What you've built}]
One more word for telling, explaining and describing at B1.
\end{grammarlens}

\subsection*{Your turn}
\begin{itemize}
\raggedright
  \item \emph{Say it:} \emph{fullō, fullōnis}
  \item \emph{Say it:} \emph{fullō, fullōnis}, once more
  \item \emph{Say it:} a sentence of your own with \emph{fullō, fullōnis}
  \item \emph{Recall:} say the Latin for a labourer, then the Latin for to work, then say \emph{fullō, fullōnis} again
\end{itemize}

\begin{tcolorbox}[breakable,colback=teal!4,colframe=teal!35!black,title={Before you move on}]
What does fullō, fullōnis mean? (\textquotedblleft{}A fuller, a cleaner of cloth\textquotedblright{}.) Say it once more.
\end{tcolorbox}
\section[aurifex, aurificis]{aurifex, aurificis — a goldsmith}

\label{lesson:LA-C377-aurifex}



\begin{quote}
Before the new one: say the Latin for to work, then the Latin for a fuller.
\end{quote}

\subsection*{You'll want to know: aurifex, aurificis}
\textbf{aurifex, aurificis} — \textquotedblleft{}a goldsmith\textquotedblright{}.

\begin{quote}
Aurifex ānulum aureum fēcit.
\end{quote}

\begin{quote}
\emph{The goldsmith made a gold ring.}
\end{quote}

\emph{aurifex, aurificis} m., from \emph{aurum}, gold, and \emph{faciō}, to make.

\begin{cousinweb}
\emph{aurifex} is built from \emph{aurum} \textquotedblleft{}gold\textquotedblright{} and \emph{faciō} \textquotedblleft{}to make\textquotedblright{} (EtymDB; Lewis \& Short): a gold-maker.
\end{cousinweb}

\begin{grammarlens}[title={What you've built}]
One more word for telling, explaining and describing at B1.
\end{grammarlens}

\subsection*{Your turn}
\begin{itemize}
\raggedright
  \item \emph{Say it:} \emph{aurifex, aurificis}
  \item \emph{Say it:} \emph{aurifex, aurificis}, once more
  \item \emph{Say it:} a sentence of your own with \emph{aurifex, aurificis}
  \item \emph{Recall:} say the Latin for to work, then the Latin for a fuller, then say \emph{aurifex, aurificis} again
\end{itemize}

\begin{tcolorbox}[breakable,colback=teal!4,colframe=teal!35!black,title={Before you move on}]
What does aurifex, aurificis mean? (\textquotedblleft{}A goldsmith\textquotedblright{}.) Say it once more.
\end{tcolorbox}
\section[lignārius, lignāriī]{lignārius, lignāriī — a carpenter, a worker in wood}

\label{lesson:LA-C377-lignarius}



\begin{quote}
Before the new one: say the Latin for a fuller, then the Latin for a goldsmith.
\end{quote}

\subsection*{You'll want to know: lignārius, lignāriī}
\textbf{lignārius, lignāriī} — \textquotedblleft{}a carpenter, a worker in wood\textquotedblright{}.

\begin{quote}
Lignārius mēnsās et sellās facit.
\end{quote}

\begin{quote}
\emph{The carpenter makes tables and chairs.}
\end{quote}

\emph{lignārius, lignāriī} m., often \emph{faber lignārius}: a smith (\emph{faber}) who works in wood, \emph{lignum}.

\begin{cousinweb}
\emph{lignārius} is built on \emph{lignum} \textquotedblleft{}wood\textquotedblright{} (EtymDB; Lewis \& Short), with the ending \emph{-ārius} of a man busy with something.
\end{cousinweb}

\begin{grammarlens}[title={What you've built}]
One more word for telling, explaining and describing at B1.
\end{grammarlens}

\subsection*{Your turn}
\begin{itemize}
\raggedright
  \item \emph{Say it:} \emph{lignārius, lignāriī}
  \item \emph{Say it:} \emph{lignārius, lignāriī}, once more
  \item \emph{Say it:} a sentence of your own with \emph{lignārius, lignāriī}
  \item \emph{Recall:} say the Latin for a fuller, then the Latin for a goldsmith, then say \emph{lignārius, lignāriī} again
\end{itemize}

\begin{tcolorbox}[breakable,colback=teal!4,colframe=teal!35!black,title={Before you move on}]
What does lignārius, lignāriī mean? (\textquotedblleft{}A carpenter, a worker in wood\textquotedblright{}.) Say it once more.
\end{tcolorbox}
\section[caementārius, caementāriī]{caementārius, caementāriī — a mason, a builder in stone}

\label{lesson:LA-C377-caementarius}



\begin{quote}
Before the new one: say the Latin for a goldsmith, then the Latin for a carpenter.
\end{quote}

\subsection*{You'll want to know: caementārius, caementāriī}
\textbf{caementārius, caementāriī} — \textquotedblleft{}a mason, a builder in stone\textquotedblright{}.

\begin{quote}
Caementāriī mūrum novum aedificant.
\end{quote}

\begin{quote}
\emph{The masons are building a new wall.}
\end{quote}

\emph{caementārius, caementāriī} m., built on \emph{caementum}, rough-cut stone set in mortar. English \emph{cement} comes from \emph{caementum}. A late word, in Jerome.

\begin{cousinweb}
\emph{caementārius} is built on \emph{caementum} \textquotedblleft{}rough-cut stone\textquotedblright{} (Lewis \& Short), which comes from \emph{caedō} \textquotedblleft{}to cut\textquotedblright{}. English \emph{cement} comes from \emph{caementum} through French.
\end{cousinweb}

\begin{grammarlens}[title={What you've built}]
One more word for telling, explaining and describing at B1.
\end{grammarlens}

\subsection*{Your turn}
\begin{itemize}
\raggedright
  \item \emph{Say it:} \emph{caementārius, caementāriī}
  \item \emph{Say it:} \emph{caementārius, caementāriī}, once more
  \item \emph{Say it:} a sentence of your own with \emph{caementārius, caementāriī}
  \item \emph{Recall:} say the Latin for a goldsmith, then the Latin for a carpenter, then say \emph{caementārius, caementāriī} again
\end{itemize}

\begin{tcolorbox}[breakable,colback=teal!4,colframe=teal!35!black,title={Before you move on}]
What does caementārius, caementāriī mean? (\textquotedblleft{}A mason, a builder in stone\textquotedblright{}.) Say it once more.
\end{tcolorbox}
\section[statuārius, statuāriī]{statuārius, statuāriī — a sculptor of statues}

\label{lesson:LA-C377-statuarius}



\begin{quote}
Before the new one: say the Latin for a carpenter, then the Latin for a mason.
\end{quote}

\subsection*{You'll want to know: statuārius, statuāriī}
\textbf{statuārius, statuāriī} — \textquotedblleft{}a sculptor of statues\textquotedblright{}.

\begin{quote}
Statuārius ex marmore deam fēcit.
\end{quote}

\begin{quote}
\emph{The sculptor made a goddess out of marble.}
\end{quote}

\emph{statuārius, statuāriī} m., built on \emph{statua}, a statue. English \emph{statuary}.

\begin{cousinweb}
\emph{statuārius} is built on \emph{statua} \textquotedblleft{}a statue\textquotedblright{} (Lewis \& Short), which comes from \emph{statuō} \textquotedblleft{}to set up\textquotedblright{}. English \emph{statuary} comes from it.
\end{cousinweb}

\begin{grammarlens}[title={What you've built}]
One more word for telling, explaining and describing at B1.
\end{grammarlens}

\subsection*{Your turn}
\begin{itemize}
\raggedright
  \item \emph{Say it:} \emph{statuārius, statuāriī}
  \item \emph{Say it:} \emph{statuārius, statuāriī}, once more
  \item \emph{Say it:} a sentence of your own with \emph{statuārius, statuāriī}
  \item \emph{Recall:} say the Latin for a carpenter, then the Latin for a mason, then say \emph{statuārius, statuāriī} again
\end{itemize}

\begin{tcolorbox}[breakable,colback=teal!4,colframe=teal!35!black,title={Before you move on}]
What does statuārius, statuāriī mean? (\textquotedblleft{}A sculptor of statues\textquotedblright{}.) Say it once more.
\end{tcolorbox}
